\documentclass[10pt,a4paper]{amsart}
\usepackage[margin=19mm]{geometry}
\usepackage{amsmath,amssymb,mathtools}
\usepackage[scr=rsfs,cal=euler]{mathalfa}
\usepackage{xfrac}
\usepackage[unicode,hidelinks,bookmarks=false]{hyperref}
\newcommand{\ie}{i.e.\ }
\newcommand{\cf}{cf.\ }
\newcommand{\resp}{respectively\ }
\newcommand{\Subsection}{\subsection}
\providecommand{\nobreakdash}{\nobreak\hbox{-}}
\setlength{\emergencystretch}{2em}
\multlinegap0pt
\usepackage{amssymb}
\usepackage{enumerate}
\usepackage{tikz}
\newtheorem{thm}{Theorem}
\newtheorem{lem}{Lemma}
\newcommand{\C}{\mathbb{C}}
\newcommand{\D}{\mathbb{D}}
\renewcommand{\Re}{\on{Re}}
\newcommand{\abs}[1]{{\left|{#1}\right|}}
\newcommand{\cG}{\mathcal{G}}
\newcommand{\fr}{\partial}
\newcommand{\inv}{^{-1}}
\newcommand{\loc}{\mathrm{loc}}
\newcommand{\lrpar}[1]{\left(#1\right)}
\newcommand{\on}[1]{\operatorname{#1}}
\newcommand{\rest}[1]{ \arrowvert_{#1}}
\newcommand{\set}[1]{{\left\{#1\right\}}}
\newcommand{\unsur}[1]{\frac{1}{#1}}
\title{On the dynamical Manin-Mumford conjecture for plane polynomial maps}
\author{Romain Dujardin, Charles Favre, Matteo Ruggiero}
\date{Source: arXiv:2312.14817v3 (CC BY 4.0)}
\begin{document}
\maketitle

\noindent\emph{Source and licence.} On the dynamical Manin-Mumford conjecture for plane polynomial maps. Version arXiv:2312.14817v3 (CC BY 4.0), distributed by arXiv under the Creative Commons Attribution 4.0 International licence, http://creativecommons.org/licenses/by/4.0/, as recorded in the arXiv licence metadata for this exact version and shown on the abstract page https://arxiv.org/abs/2312.14817v3. Full licence notice: \texttt{licenses/arxiv-2312.14817-CC-BY-4.0.txt}.\par\smallskip

\noindent\textit{Editorial scope.} Counted unit: the article's thm thm:graph\_transform\_parabolic (source0.tex lines 982--992). The article's complete original proof of it is delivered with it (source0.tex lines 1049--1100, one \texttt{proof} environment of 52 lines, which the article places away from the statement). No mathematics is written, repaired or re-derived: the statement and the proof are the article's own lines, in the article's own notation, and no retained line is substituted or re-flowed.  Retained uncounted as the source-local context the proof needs: lines 962--964.  Nothing was omitted from any retained span, so the package carries no omission record.  No retained line contains a citation, so the wrapper carries no bibliography and no bibliography entry.  No retained line contains a figure environment, an image-inclusion command or drawing code, so no image is delivered and the package declares no asset dependency.  Editorial formatting only, in addition: an AMS \texttt{amsart} preamble that restates the article's own declarations and macros used by the retained lines, namely the environments \texttt{thm}, \texttt{lem} on separate counters, together with the standard packages those lines need.  The retained environments print in this document as Theorem 1, Lemma 1 in order of appearance, whereas the article prints the same items with the numbering of its own full document.  The complete native source of the article as distributed in the arXiv e-print for this exact version is bundled unchanged as \texttt{sources/151979/source0.tex}.
\begin{equation}\label{eq:f_parabolic} 
f(x,y) =  \left(x + g(x,y) , y^d(1+h(x,y))\right)~,
\end{equation}

\begin{thm}\label{thm:graph_transform_parabolic} 
Let $f\colon(\C^2, 0) \to (\C^2, 0)$ 
  be of the form~\eqref{eq:f_parabolic} 
  with $h(0)=0$, $g(x,0)\not\equiv 0$ and $g(x,y) = \mathcal{O}(|(x,y)|^2)$. 
Let $U$ be any repelling petal of   
$f\rest{\set{y=0}}$ and consider a germ $V$ of analytic curve transverse 
to $\set{y=0}$, intersecting $U$.
 
Then there exists $r>0$  depending only on $f$, such that  for  large enough $n$, 
the analytic sets $f^{-n}(V)$ are   vertical graphs in $\D_r^2$  
converging  to $W^{\rm ss}_\loc(0)$ in the $C^1$ topology. \end{thm}

\begin{proof}[Proof of Theorem~\ref{thm:graph_transform_parabolic}]
By the previous lemma, we may suppose that 
\[f(x,y) = \left(x+x^{k+1}+ x^{2k+1} g(x,y), y^d (1+ x h(x,y))\right)\] with $g$ and $h$ holomorphic near the origin. 
Note that 
$f\rest{y=0}$ has a repelling petal along the positive real axis. 
Fix $r>0$ such that
$f$ is holomorphic and injective on a neighborhood of $\overline\D_r^2$. 
 Reducing and rotating the petal if necessary, we may assume that 
\[U:=\set{x:\ \arg(x)\in \lrpar{-\frac{\pi}{4k}, \frac{\pi}{4k}} , \ \abs{x}<r}.\] 
The holomorphic map $ z =   \lrpar{kx^k}\inv$ is univalent on 
$U\times \D_r$, and  takes its values in 
 \[\Omega_R := \set{z: \ \arg(z)\in  \lrpar{-\frac{\pi}{4}, \frac{\pi}{4}}, \ \abs{z}>R},\] where 
 $R=(kr^k)^{-1}$.
  The expression of 
 $f$ in the coordinates $(z,y)$ is of the form
\begin{equation}\label{eq:f_parabolic3}
f(z,y)  = \lrpar{z-1+ \unsur{z} a(z,y),  y^d \lrpar{1+  \unsur{z^{1/k}} b(z,y)}}. 
\end{equation}
so that $f$ is now defined in $\Omega_R\times \D_r$. 
 Fix    $M>0$ such that 
 \begin{equation}\label{eq:estimates}
 \abs{a}, \ \abs{b}, \abs{\frac{\fr a}{\fr z}}, \abs{\frac{\fr a}{\fr y}}, \ \abs{\frac{\fr b}{\fr z}}, \ \abs{\frac{\fr b}{\fr y}}
 \leq M \text{ on }\Omega_R\times \D_r, 
 \end{equation} 
and reduce $r$ if necessary so that $r<\frac1{10 d}$ and $MR^{-1/k}\leq \frac{1}{100}$.

 For any $\rho<r$ and $\sigma >0$ we let 
 \[\cG(\rho, \sigma) =\{\varphi \colon \D_\rho \to \Omega_R \text{ holomorphic  s.t. } \sup_{\D_\rho} 
 \abs{\varphi'} \le \sigma\}~.\]
\begin{lem}\label{lem:vertest}
Suppose that $\sigma \rho < \frac1{100 d}$.  
 For any vertical graph $\Gamma$ 
determined by $\varphi\in \cG(\rho, \sigma)$, then $f\inv \Gamma$ is 
a vertical graph determined by a function   $\psi \in \cG(\rho_1, \frac1{10})$
where
$\rho_1  = \min\lrpar{ (\rho/2)^{1/d}, r}$ and  $\Re(\psi) \geq \Re(\varphi(0)) + 9/10$.
\end{lem}
 
Assuming this result for the moment, 
 let us  conclude the proof of the theorem.  
Let  $V$ be any germ of curve intersecting transversely   $\set{y=0}$ 
 at $(x_0, 0)\in U$. Then $V$ is a graph of slope $\sigma$ over some  
  disk $\D_\rho$ in the second coordinate, for some
 $\rho>0$. Reduce $\rho$ by trimming $V$ if necessary so that $\sigma\rho< 1/100d$.
Since $r/10 < 1/100d$, the  previous lemma implies that 
 we can define inductively a sequence of vertical graphs
$V = V_0$,  $V_n :=  f\inv V_{n-1}$, where $V_n$ is defined by a holomorphic function $z = \varphi_n(y)$ of uniformly bounded slope %$\le \frac1{10}$
over $\D_{\rho_n}$, and furthermore  $\rho_n = r$ for  large enough   $n$.   
Moreover, we have $\Re(\varphi_n) \geq \Re(\varphi(0)) + 9n/10$, hence, coming back to the 
$(x,y)$ coordinates, we see that  $V_n = \{ x= (k\varphi_n(y))^{-1/k}\}$
converges in the $C^1$-topology to the curve $W^{\rm ss}_\loc(0)= \{x=0\}$.
\end{proof}

\end{document}
